\section{Transisi (transition)}

\textbf{Transition} mengubah nilai properti CSS secara halus dalam jangka waktu tertentu, menciptakan efek animasi sederhana yang dipicu oleh perubahan state (misalnya \code{:hover}, \code{:focus}, penambahan class via JavaScript). Transition membuat antarmuka terasa lebih responsif dan profesional tanpa memerlukan kode JavaScript \cite{mdnCss}.

\begin{table}[htbp]
\centering
\begin{tabular}{|P{5cm}|P{3.5cm}|P{3.5cm}|}
\hline
\textbf{Properti} & \textbf{Fungsi} & \textbf{Contoh Nilai} \\
\hline
\code{transition-property} & Properti yang dianimasikan & \code{background}, \code{all} \\
\hline
\code{transition-duration} & Durasi animasi & \code{0.3s}, \code{200ms} \\
\hline
\code{transition-timing-function} & Kurva kecepatan & \code{ease}, \code{linear} \\
\hline
\code{transition-delay} & Jeda sebelum mulai & \code{0s}, \code{0.1s} \\
\hline
\code{transition} (shorthand) & Gabungan semua & \code{all 0.3s ease 0s} \\
\hline
\end{tabular}
\caption{Properti transition CSS}
\end{table}

Hanya properti yang memiliki nilai \textbf{interpolable} (dapat dihitung perantaranya) yang dapat ditransisi---termasuk warna, ukuran, opacity, transform, dan padding. Properti diskret seperti \code{display} dan \code{visibility} tidak dapat ditransisi (kecuali \code{visibility} yang memiliki perilaku khusus). Untuk performa optimal, animasikan hanya \code{transform} dan \code{opacity}---keduanya dioptimalkan oleh GPU browser \cite{webdevCss}.

\begin{konsep}
\textbf{Performa Animasi.} Properti yang mempengaruhi layout (width, height, margin, padding) memicu \textit{reflow}---proses mahal di mana browser menghitung ulang posisi semua elemen. Properti seperti \code{transform} dan \code{opacity} hanya memicu \textit{compositing} (proses ringan). Selalu prioritaskan animasi dengan transform/opacity untuk pengalaman yang mulus, terutama di perangkat mobile \cite{cssTricks}.
\end{konsep}

